\documentclass[12pt,a4paper]{article}

\usepackage[margin=1in]{geometry}
\usepackage{graphicx}
\usepackage{booktabs}
\usepackage{amsmath}
\usepackage{hyperref}
\usepackage{float}
\usepackage{caption}
\usepackage{subcaption}
\usepackage{array}
\usepackage{xcolor}
\usepackage{titlesec}
\usepackage{parskip}

\hypersetup{
    colorlinks=true,
    linkcolor=blue,
    urlcolor=blue,
}

\title{
    \Large \textbf{DLOps Assignment 5} \\[0.5em]
    \large Question 2: Adversarial Attacks using IBM Adversarial Robustness Toolkit (ART)
}
\author{Roll No: M25CSA018}
\date{}

\begin{document}

\maketitle

\tableofcontents
\newpage

% ─────────────────────────────────────────────────────────────────────────────
\section{Introduction}

This report covers the implementation and evaluation of adversarial attacks and
defenses on image classification models using the IBM Adversarial Robustness
Toolkit (ART). The assignment is divided into two parts:

\begin{itemize}
    \item \textbf{Part (i):} Implementing the Fast Gradient Sign Method (FGSM)
    attack from scratch and via IBM ART on a ResNet18 model trained on CIFAR-10,
    and analyzing the impact of perturbation strength on model accuracy.

    \item \textbf{Part (ii):} Designing and training deep learning-based
    adversarial detectors using ResNet34 to classify clean vs.\ adversarial
    images generated by PGD and BIM attacks.
\end{itemize}

All experiments were run inside a Docker container using Python and PyTorch.
Experiment tracking was done via Weights \& Biases (WandB).

\textbf{WandB Project:}
\url{https://wandb.ai/m25csa018-iit-j/dlops-q2-adversarial}

\textbf{GitHub Branch (Assignment-5):}
\url{https://github.com/S-Ayush/Sample_repo/tree/Assignment-5}

% ─────────────────────────────────────────────────────────────────────────────
\section{Part (i): FGSM Attack — From Scratch vs IBM ART}

\subsection{Dataset and Model}

\textbf{Dataset:} CIFAR-10 — 60,000 color images ($32 \times 32$) across
10 classes (airplane, automobile, bird, cat, deer, dog, frog, horse, ship,
truck). 50,000 training and 10,000 test images.

\textbf{Model:} ResNet18 trained from scratch (no pretrained weights).
The standard ImageNet ResNet18 was adapted for CIFAR-10 by:
\begin{itemize}
    \item Replacing the $7\times7$ stride-2 convolution with a $3\times3$
    stride-1 convolution to preserve spatial resolution on $32\times32$ inputs.
    \item Removing the MaxPool layer to avoid aggressive downsampling.
    \item Replacing the final fully connected layer with a 10-class output head.
\end{itemize}

\textbf{Training Configuration:}
\begin{center}
\begin{tabular}{ll}
\toprule
\textbf{Hyperparameter} & \textbf{Value} \\
\midrule
Optimizer       & SGD with Nesterov momentum \\
Learning Rate   & 0.1 (CosineAnnealingLR over 50 epochs) \\
Momentum        & 0.9 \\
Weight Decay    & $5 \times 10^{-4}$ \\
Batch Size      & 128 \\
Epochs          & 50 \\
Augmentation    & RandomCrop(32, pad=4), RandomHorizontalFlip \\
\bottomrule
\end{tabular}
\end{center}

\subsection{Classifier Results}

\begin{center}
\begin{tabular}{ll}
\toprule
\textbf{Metric} & \textbf{Value} \\
\midrule
Best Validation Accuracy & \textbf{94.31\%} \\
Final Training Accuracy  & 99.88\% \\
Requirement              & $\geq 72\%$ \textcolor{green}{$\checkmark$} \\
\bottomrule
\end{tabular}
\end{center}

The model comfortably exceeds the 72\% accuracy requirement.

\subsection{FGSM Attack: From Scratch}

The Fast Gradient Sign Method (FGSM) \cite{goodfellow2014} generates adversarial
examples by taking a single gradient step in the direction that maximizes the
classification loss:

\begin{equation}
    x_{\text{adv}} = \text{clip}\left( x + \varepsilon \cdot
    \text{sign}\left( \nabla_x \mathcal{L}(\theta, x, y) \right),\ 0,\ 1 \right)
\end{equation}

where $x$ is the original image, $y$ is the true label, $\mathcal{L}$ is the
cross-entropy loss, $\theta$ are the model weights, and $\varepsilon$ is the
perturbation budget.

\textbf{Implementation steps:}
\begin{enumerate}
    \item Enable gradient tracking on the input tensor (\texttt{requires\_grad = True})
    \item Forward pass to compute the classification loss
    \item Backward pass to compute $\nabla_x \mathcal{L}$
    \item Apply the sign of the gradient scaled by $\varepsilon$
    \item Clip pixel values to $[0, 1]$
\end{enumerate}

\subsection{FGSM Attack: IBM ART}

IBM ART's \texttt{FastGradientMethod} was used with a
\texttt{PyTorchClassifier} wrapper. ART internally handles normalization via its
preprocessing pipeline, so raw $[0, 1]$ images were passed directly.

\subsection{Attack Results ($\varepsilon = 0.03$)}

\begin{center}
\begin{tabular}{lrrr}
\toprule
\textbf{Attack} & \textbf{Clean Acc} & \textbf{Adv Acc} & \textbf{Drop} \\
\midrule
No Attack       & 94.31\% & 94.31\% &  0.00\% \\
FGSM (Scratch)  & 94.31\% & 32.77\% & 61.54\% \\
FGSM (ART)      & 94.31\% & 36.07\% & 58.24\% \\
PGD  (ART)      & 94.31\% &  3.45\% & \textbf{90.86\%} \\
BIM  (ART)      & 94.31\% &  3.90\% & \textbf{90.41\%} \\
\bottomrule
\end{tabular}
\end{center}

\subsection{Visual Comparison}

\begin{figure}[H]
    \centering
    \includegraphics[width=\textwidth]{fgsm_comparison.png}
    \caption{Original vs FGSM (Scratch) vs FGSM (ART) — 10 sample images.
    Perturbations are invisible to the human eye yet cause dramatic accuracy drops.}
    \label{fig:fgsm_comparison}
\end{figure}

\begin{figure}[H]
    \centering
    \begin{subfigure}[b]{0.48\textwidth}
        \includegraphics[width=\textwidth]{pgd_comparison.png}
        \caption{Original vs PGD adversarial images}
    \end{subfigure}
    \hfill
    \begin{subfigure}[b]{0.48\textwidth}
        \includegraphics[width=\textwidth]{bim_comparison.png}
        \caption{Original vs BIM adversarial images}
    \end{subfigure}
    \caption{PGD and BIM attack comparisons. Both reduce accuracy to below 4\%.}
    \label{fig:pgd_bim_comparison}
\end{figure}

\subsection{Perturbation Strength vs Accuracy Drop}

The epsilon sweep evaluates FGSM (ART) at multiple perturbation budgets:

\begin{center}
\begin{tabular}{rrrr}
\toprule
\textbf{$\varepsilon$} & \textbf{Clean Acc} & \textbf{Adv Acc} & \textbf{Drop} \\
\midrule
0.010 & 94.31\% & 47.43\% & 46.88\% \\
0.020 & 94.31\% & 39.90\% & 54.41\% \\
0.030 & 94.31\% & 36.07\% & 58.24\% \\
0.050 & 94.31\% & 28.65\% & 65.66\% \\
0.100 & 94.31\% & 14.78\% & 79.53\% \\
\bottomrule
\end{tabular}
\end{center}

\begin{figure}[H]
    \centering
    \includegraphics[width=\textwidth]{epsilon_sweep.png}
    \caption{Accuracy vs epsilon (left) and accuracy drop vs epsilon (right).
    As $\varepsilon$ increases, adversarial accuracy decreases monotonically.}
    \label{fig:epsilon_sweep}
\end{figure}

\subsection{Analysis}

\textbf{FGSM Scratch vs ART:} Both methods produce similar accuracy drops
(61.54\% vs 58.24\%), confirming the scratch implementation is correct.
The small difference ($\approx 3\%$) is due to minor differences in
normalization handling.

\textbf{PGD vs BIM vs FGSM:} PGD and BIM are substantially more powerful
than FGSM, reducing accuracy to $\approx 3\%$ vs $\approx 33\%$ for FGSM.
This is because PGD and BIM apply 40 iterative gradient steps, each one
maximizing the loss further. PGD additionally uses a random start, making
it marginally stronger than BIM.

\textbf{Epsilon effect:} Accuracy drops monotonically as $\varepsilon$
increases. At $\varepsilon = 0.1$, accuracy falls to 14.78\%, yet the
perturbation becomes slightly visible to the human eye. At $\varepsilon
= 0.01$, the perturbation is completely invisible but still causes a
46.88\% accuracy drop — demonstrating the vulnerability of deep neural
networks to imperceptible input modifications.

% ─────────────────────────────────────────────────────────────────────────────
\section{Part (ii): Adversarial Detection}

\subsection{Task Overview}

Two binary classifiers (detectors) were trained using ResNet34 to distinguish
between clean and adversarial images:

\begin{itemize}
    \item \textbf{Detector A:} Trained on clean images + PGD adversarial images
    \item \textbf{Detector B:} Trained on clean images + BIM adversarial images
\end{itemize}

Labels: $0 = \text{clean}$, $1 = \text{adversarial}$.

\subsection{Adversarial Image Generation}

\textbf{PGD (Projected Gradient Descent):}
\begin{equation}
    x^{t+1} = \text{Clip}_{\varepsilon}\left( x^t + \alpha \cdot
    \text{sign}\left(\nabla_x \mathcal{L}(\theta, x^t, y)\right) \right)
\end{equation}
Parameters: $\varepsilon = 0.03$, $\alpha = 0.007$, iterations $= 40$,
random initialization enabled.

\textbf{BIM (Basic Iterative Method):}
Same as PGD but without random initialization — deterministic iterative FGSM.

\subsection{Detector Architecture}

ResNet34 adapted for CIFAR-10 ($32 \times 32$) with the same modifications
as ResNet18: $3\times3$ stride-1 first convolution, no MaxPool, and a
2-class output head.

\textbf{Training Configuration:}
\begin{center}
\begin{tabular}{ll}
\toprule
\textbf{Hyperparameter} & \textbf{Value} \\
\midrule
Optimizer       & Adam (lr=$10^{-3}$, weight decay=$10^{-4}$) \\
Scheduler       & StepLR (step=5, $\gamma=0.5$) \\
Epochs          & 15 \\
Batch Size      & 128 \\
Dataset         & 10,000 clean + 10,000 adversarial (80/20 split) \\
\bottomrule
\end{tabular}
\end{center}

\subsection{Detector Results}

\begin{center}
\begin{tabular}{lrrrr}
\toprule
\textbf{Detector} & \textbf{Overall Acc} & \textbf{Clean Det} & \textbf{Adv Det} & \textbf{Status} \\
\midrule
PGD Detector & \textbf{99.97\%} & 100.00\% & 99.95\% & \textcolor{green}{PASS} \\
BIM Detector & \textbf{99.97\%} & 99.95\%  & 100.00\% & \textcolor{green}{PASS} \\
Requirement  & $\geq 70\%$      & —        & —        & \\
\bottomrule
\end{tabular}
\end{center}

\begin{figure}[H]
    \centering
    \includegraphics[width=0.75\textwidth]{detector_comparison.png}
    \caption{PGD vs BIM detector comparison. Both achieve near-perfect
    detection accuracy, well above the 70\% requirement.}
    \label{fig:detector_comparison}
\end{figure}

\subsection{Analysis}

Both detectors achieved $\approx 99.97\%$ overall accuracy, far exceeding
the 70\% requirement. This indicates that PGD and BIM adversarial images,
despite being imperceptible to humans, leave a strong statistical signature
that a classifier can learn to detect.

\textbf{PGD vs BIM detection:} The detection accuracy is nearly identical
for both attacks (99.97\% each). This is expected because BIM is a special
case of PGD (without random initialization), so the adversarial perturbation
patterns they produce are structurally similar, making them equally
detectable.

\textbf{Why detection is so accurate:} Adversarial perturbations generated
by gradient-based attacks tend to add structured high-frequency noise
patterns that differ systematically from natural image noise. The ResNet34
detector learns to recognize these patterns. The near-perfect score also
suggests the detector may be partially learning attack-specific artifacts
rather than general adversarial patterns — cross-attack generalization
would require training on a mixture of attack types.

% ─────────────────────────────────────────────────────────────────────────────
\section{Conclusion}

\begin{itemize}
    \item ResNet18 trained from scratch achieved \textbf{94.31\%} accuracy on
    CIFAR-10, well above the 72\% requirement.

    \item Adversarial attacks successfully degraded model performance:
    FGSM reduced accuracy to $\approx 33\%$, while PGD and BIM reduced it
    to $\approx 3\%$ — demonstrating that iterative attacks are substantially
    more powerful.

    \item The FGSM scratch implementation matched IBM ART's implementation
    closely ($\Delta \approx 3\%$), validating the manual implementation.

    \item Both adversarial detectors (ResNet34) achieved $\approx 99.97\%$
    detection accuracy on PGD and BIM adversarial images, far exceeding
    the 70\% requirement.
\end{itemize}

% ─────────────────────────────────────────────────────────────────────────────
\begin{thebibliography}{9}

\bibitem{goodfellow2014}
Goodfellow, I.J., Shlens, J., \& Szegedy, C. (2014).
\textit{Explaining and harnessing adversarial examples}.
arXiv:1412.6572.

\bibitem{madry2017}
Madry, A., Makelov, A., Schmidt, L., Tsipras, D., \& Vladu, A. (2017).
\textit{Towards deep learning models resistant to adversarial attacks}.
arXiv:1706.06083.

\bibitem{kurakin2016}
Kurakin, A., Goodfellow, I., \& Bengio, S. (2016).
\textit{Adversarial examples in the physical world}.
arXiv:1607.02533.

\bibitem{art2018}
Nicolae, M.I., et al. (2018).
\textit{Adversarial Robustness Toolbox v1.0.0}.
arXiv:1807.01069.

\end{thebibliography}

\end{document}
